% Options for packages loaded elsewhere
\PassOptionsToPackage{unicode}{hyperref}
\PassOptionsToPackage{hyphens}{url}
\documentclass[
  12pt,
]{ctexart}
\usepackage{xcolor}
\usepackage{amsmath,amssymb}
\setcounter{secnumdepth}{-\maxdimen} % remove section numbering
\usepackage{iftex}
\ifPDFTeX
  \usepackage[T1]{fontenc}
  \usepackage[utf8]{inputenc}
  \usepackage{textcomp} % provide euro and other symbols
\else % if luatex or xetex
  \usepackage{unicode-math} % this also loads fontspec
  \defaultfontfeatures{Scale=MatchLowercase}
  \defaultfontfeatures[\rmfamily]{Ligatures=TeX,Scale=1}
\fi
\usepackage{lmodern}
\ifPDFTeX\else
  % xetex/luatex font selection
\fi
% Use upquote if available, for straight quotes in verbatim environments
\IfFileExists{upquote.sty}{\usepackage{upquote}}{}
\IfFileExists{microtype.sty}{% use microtype if available
  \usepackage[]{microtype}
  \UseMicrotypeSet[protrusion]{basicmath} % disable protrusion for tt fonts
}{}
\makeatletter
\@ifundefined{KOMAClassName}{% if non-KOMA class
  \IfFileExists{parskip.sty}{%
    \usepackage{parskip}
  }{% else
    \setlength{\parindent}{0pt}
    \setlength{\parskip}{6pt plus 2pt minus 1pt}}
}{% if KOMA class
  \KOMAoptions{parskip=half}}
\makeatother
\usepackage{longtable,booktabs,array}
\usepackage{calc} % for calculating minipage widths
% Correct order of tables after \paragraph or \subparagraph
\usepackage{etoolbox}
\makeatletter
\patchcmd\longtable{\par}{\if@noskipsec\mbox{}\fi\par}{}{}
\makeatother
% Allow footnotes in longtable head/foot
\IfFileExists{footnotehyper.sty}{\usepackage{footnotehyper}}{\usepackage{footnote}}
\makesavenoteenv{longtable}
\setlength{\emergencystretch}{3em} % prevent overfull lines
\providecommand{\tightlist}{%
  \setlength{\itemsep}{0pt}\setlength{\parskip}{0pt}}
\usepackage{bookmark}
\IfFileExists{xurl.sty}{\usepackage{xurl}}{} % add URL line breaks if available
\urlstyle{same}
\hypersetup{
  hidelinks,
  pdfcreator={LaTeX via pandoc}}

\author{SCX}
\date{2026}

\begin{document}

\section{PAC-Bayes Generalization Bound for the SCX Noise
Detector}\label{pac-bayes-generalization-bound-for-the-scx-noise-detector}

\begin{quote}
\textbf{Paper 3 Upgrade Item 1} \textbf{Status}: Draft --- proof
complete modulo standard PAC-Bayes machinery \textbf{Target location}:
Paper 3 Section 4.2 (Theorem 5) \textbf{Date}: 2026-06-27
\end{quote}

\begin{center}\rule{0.5\linewidth}{0.5pt}\end{center}

\subsection{1. Motivation}\label{motivation}

\textbf{Current limitation of Theorem 1}: Theorem 1 bounds the true F1
assuming knowledge of the state-level expert error rates
\(\{\mu_s\}_{s \in \mathcal{S}}\):

\[\text{F1} \geq 1 - \frac{1}{\eta} \sum_{s \in \mathcal{S}} \rho_s \cdot \exp\!\bigl(-2M\Delta_s^2\bigr)\]

where
\(\Delta_s = \min(\theta - \mu_s,\; 1 - C_{\text{bal}} \cdot \mu_s/(K-1) - \theta)\).
In practice, \(\mu_s\) is \textbf{unknown} and must be estimated from a
finite clean validation set. The threshold \(\theta\) is then chosen
based on the estimates \(\{\hat{\mu}_s\}\).

\textbf{The problem}: When we plug \(\hat{\mu}_s\) into Theorem 1, the
resulting bound is an \emph{estimate} of the true F1. We need a bound
that holds with high probability over the randomness of the validation
set, accounting for the complexity of the threshold selection procedure.

\textbf{The solution}: PAC-Bayes theory (McAllester, 1999; Catoni, 2007)
provides a principled way to bound the gap between empirical and true
performance when the hypothesis (here: the detection threshold) is
chosen in a data-dependent manner. The penalty depends on the KL
divergence between a data-dependent posterior and a fixed prior,
measuring the ``information cost'' of adapting the threshold to the
data.

\textbf{Value}: This makes the F1 bound fully empirical --- no oracle
quantities required.

\begin{center}\rule{0.5\linewidth}{0.5pt}\end{center}

\subsection{2. Setup and Notation}\label{setup-and-notation}

\subsubsection{2.1 The Noise Detector as a
Classifier}\label{the-noise-detector-as-a-classifier}

The SCX noise detector \(h_\theta: \mathcal{X} \to \{0,1\}\) is a binary
classifier:

\[h_\theta(x, y, \{f_m\}_{m=1}^M) = \mathbf{1}\{C(x) > \theta\}, \quad
C(x) = \frac{1}{M} \sum_{m=1}^M \mathbf{1}\{\ell(f_m(x), y) > \tau\}\]

where output \(1\) means ``predicted noise'' and \(0\) means ``predicted
clean.''

\subsubsection{2.2 Loss Function}\label{loss-function}

Let \(z \in \{0,1\}\) be the true noise indicator (\(z = 1\) means the
label is corrupted). Define the \textbf{detection loss}:

\[\ell_\theta(x, y, z) = \mathbf{1}\{h_\theta(x) \neq z\}\]

This is the 0-1 classification loss for the noise detection task. Its
expectation is the \textbf{detection error rate}:

\[R(\theta) = \mathbb{E}[\ell_\theta(X, Y, Z)] = \eta \cdot \text{FNR}(\theta) + (1-\eta) \cdot \text{FPR}(\theta)\]

\subsubsection{2.3 Relation to F1 Score}\label{relation-to-f1-score}

The F1 score is connected to \(R(\theta)\) through the precision-recall
tradeoff. Specifically:

\[\text{F1}(\theta) = \frac{2\eta(1 - \text{FNR}(\theta))}{2\eta(1 - \text{FNR}(\theta)) + (1-\eta)\text{FPR}(\theta) + \eta \cdot \text{FNR}(\theta)}\]

For small \(\text{FPR}\) and \(\text{FNR}\) (the regime of interest in
Theorem 1), we have the approximation:

\[1 - \text{F1}(\theta) \approx \frac{(1-\eta)\text{FPR}(\theta) + \eta \cdot \text{FNR}(\theta)}{\eta} = \frac{R(\theta)}{\eta}\]

More precisely, we have the following inequality (derived in Lemma 4
below):

\[1 - \text{F1}(\theta) \leq \frac{R(\theta)}{\eta \cdot (1 - R(\theta))}\]

For \(R(\theta) \leq 1/2\) (which holds under Theorem 1 conditions),
this simplifies to:

\[1 - \text{F1}(\theta) \leq \frac{2R(\theta)}{\eta}\]

\subsubsection{2.4 Validation Set}\label{validation-set}

We have a labeled validation set:

\[V = \{(x_i, y_i, z_i)\}_{i=1}^N \sim \mathcal{D}^N\]

where \(z_i = \mathbf{1}\{y_i \neq y_i^*\}\) is the true noise
indicator. This validation set may be: - \textbf{Synthetic}: Take clean
data, inject known noise at rate \(\eta\) - \textbf{Natural}: A held-out
set with ground-truth labels (e.g., from a domain expert)

The empirical detection error is:

\[\hat{R}_V(\theta) = \frac{1}{N} \sum_{i=1}^N \mathbf{1}\{h_\theta(x_i, y_i) \neq z_i\}\]

\subsubsection{2.5 PAC-Bayes Framework}\label{pac-bayes-framework}

Let \(\Theta \subseteq [0,1]\) be the hypothesis class of thresholds.
Let \(P\) be a \textbf{prior distribution} over \(\Theta\) (chosen
before seeing \(V\)). Let \(Q\) be a \textbf{posterior distribution}
over \(\Theta\) (possibly data-dependent, e.g., concentrated around an
empirically optimal threshold).

The expected risk under \(Q\) is:

\[R(Q) = \mathbb{E}_{\theta \sim Q}[R(\theta)], \quad
\hat{R}_V(Q) = \mathbb{E}_{\theta \sim Q}[\hat{R}_V(\theta)]\]

\begin{center}\rule{0.5\linewidth}{0.5pt}\end{center}

\subsection{3. Main Theorem}\label{main-theorem}

\subsubsection{Theorem 5 (PAC-Bayes SCX
Generalization)}\label{theorem-5-pac-bayes-scx-generalization}

Let the setup of Section 2 hold. Let \(P\) be any prior distribution
over \(\Theta\) that does not depend on \(V\). Assume the loss
\(\ell_\theta \in [0,1]\) (bounded). Then for any \(\delta \in (0,1)\),
with probability at least \(1-\delta\) over draws of
\(V \sim \mathcal{D}^N\), for \textbf{all} distributions \(Q\) over
\(\Theta\) (possibly data-dependent):

\[R(Q) \leq \hat{R}_V(Q) + \sqrt{\frac{\text{KL}(Q \parallel P) + \log\frac{2\sqrt{N}}{\delta}}{2N}}\]

where \(\text{KL}(Q \parallel P) = \int \log\frac{dQ}{dP}\, dQ\) is the
Kullback-Leibler divergence.

\subsubsection{Corollary 5.1 (F1 Form)}\label{corollary-5.1-f1-form}

Under the same conditions, if \(R(\theta) \leq 1/2\) for all
\(\theta \in \Theta\) (which holds when \(\Delta_s > 0\), i.e., under
B3), then with probability at least \(1-\delta\):

\[1 - \text{F1}(Q) \leq \frac{2}{\eta}\left[\hat{R}_V(Q) + \sqrt{\frac{\text{KL}(Q \parallel P) + \log\frac{2\sqrt{N}}{\delta}}{2N}}\right]\]

Equivalently:

\[\text{F1}(Q) \geq \widehat{\text{F1}}(Q) - \frac{2}{\eta}\sqrt{\frac{\text{KL}(Q \parallel P) + \log\frac{2\sqrt{N}}{\delta}}{2N}}\]

where
\(\widehat{\text{F1}}(Q) = \mathbb{E}_{\theta\sim Q}[\widehat{\text{F1}}_V(\theta)]\)
is the expected empirical F1 on the validation set.

\subsubsection{Corollary 5.2 (Threshold Learning
Bound)}\label{corollary-5.2-threshold-learning-bound}

For the special case where \(Q = \delta_{\hat{\theta}}\) is a Dirac at
the empirically optimal threshold
\(\hat{\theta} = \arg\min_\theta \hat{R}_V(\theta)\), and \(P\) is a
uniform prior over \(\Theta\), we have
\(\text{KL}(\delta_{\hat{\theta}} \parallel P) = \log|\Theta|\) (where
\(|\Theta|\) is the number of candidate thresholds in a discretized
grid). Then:

\[\text{F1}(\hat{\theta}) \geq \widehat{\text{F1}}_V(\hat{\theta}) - \frac{2}{\eta}\sqrt{\frac{\log|\Theta| + \log\frac{2\sqrt{N}}{\delta}}{2N}}\]

This shows that the penalty for threshold selection grows only
logarithmically in the number of candidate thresholds --- consistent
with the Vapnik-Chervonenkis theory for threshold classifiers.

\begin{center}\rule{0.5\linewidth}{0.5pt}\end{center}

\subsection{4. Proof}\label{proof}

\subsubsection{4.1 Lemma 4 (F1-Risk
Relation)}\label{lemma-4-f1-risk-relation}

\textbf{Lemma 4.} For any noise detector \(h\) with detection error rate
\(R = \mathbb{P}(h(X) \neq Z)\) and F1 score \(\text{F1}\):

\[1 - \text{F1} \leq \frac{R}{\eta \cdot (1 - R)}\]

If \(R \leq 1/2\), then:

\[1 - \text{F1} \leq \frac{2R}{\eta}\]

\textbf{Proof.} Recall the definitions:

\[\begin{aligned}
\text{TPR} &= \mathbb{P}(h = 1 \mid Z = 1), \quad
\text{FPR} = \mathbb{P}(h = 1 \mid Z = 0), \\
\text{FNR} &= 1 - \text{TPR} = \mathbb{P}(h = 0 \mid Z = 1)
\end{aligned}\]

The detection error rate is:

\[R = \eta \cdot \text{FNR} + (1-\eta) \cdot \text{FPR}\]

F1 is:

\[\text{F1} = \frac{2\eta \cdot \text{TPR}}{2\eta \cdot \text{TPR} + (1-\eta) \cdot \text{FPR} + \eta \cdot \text{FNR}}\]

Substituting \(\text{TPR} = 1 - \text{FNR}\):

\[\begin{aligned}
\text{F1} &= \frac{2\eta(1 - \text{FNR})}{2\eta(1 - \text{FNR}) + (1-\eta)\text{FPR} + \eta \cdot \text{FNR}} \\
&= \frac{2\eta - 2\eta\cdot\text{FNR}}{2\eta - 2\eta\cdot\text{FNR} + (1-\eta)\text{FPR} + \eta\cdot\text{FNR}} \\
&= \frac{2\eta - 2\eta\cdot\text{FNR}}{2\eta - \eta\cdot\text{FNR} + (1-\eta)\text{FPR}}
\end{aligned}\]

Therefore:

\[\begin{aligned}
1 - \text{F1} &= 1 - \frac{2\eta - 2\eta\cdot\text{FNR}}{2\eta - \eta\cdot\text{FNR} + (1-\eta)\text{FPR}} \\
&= \frac{(1-\eta)\text{FPR} + \eta\cdot\text{FNR}}{2\eta - \eta\cdot\text{FNR} + (1-\eta)\text{FPR}} \\
&= \frac{R}{2\eta - \eta\cdot\text{FNR} + (1-\eta)\text{FPR}} \\
&\leq \frac{R}{\eta(2 - \text{FNR})} \quad (\text{since } (1-\eta)\text{FPR} \geq 0) \\
&\leq \frac{R}{\eta} \quad (\text{since } \text{FNR} \leq 1 \Rightarrow 2 - \text{FNR} \geq 1)
\end{aligned}\]

For the refined bound, note that
\(R = \eta \cdot \text{FNR} + (1-\eta) \cdot \text{FPR}\) and
\(2\eta - \eta\cdot\text{FNR} + (1-\eta)\text{FPR} = \eta(2 - \text{FNR}) + (1-\eta)\text{FPR}\).
Since \(2 - \text{FNR} \geq 1\), the denominator is at least \(\eta\).
Moreover, when \(R \leq 1/2\), we have \(1 - R \geq 1/2\), so:

\[1 - \text{F1} = \frac{R}{2\eta - \eta\cdot\text{FNR} + (1-\eta)\text{FPR}} \leq \frac{R}{2\eta - \eta\cdot\text{FNR}} \leq \frac{R}{\eta}\]

The tighter bound \(1 - \text{F1} \leq R/(\eta \cdot (1 - R))\) follows
from noting that:

\[\begin{aligned}
2\eta - \eta\cdot\text{FNR} + (1-\eta)\text{FPR}
&= \eta(1 + \text{TPR}) + (1-\eta)\text{FPR} \\
&\geq \eta(1 + \text{TPR}) \geq \eta(1 + 1 - \text{FNR}) = \eta(2 - \text{FNR})
\end{aligned}\]

And since \(R \geq \eta \cdot \text{FNR}\):

\[2 - \text{FNR} \geq 2 - \frac{R}{\eta} \geq 1 - R\]

when \(R \leq \eta \leq 1/2\). In general,
\(1 - \text{F1} = R/(\eta \cdot \text{TPR} + (1-\eta)(1 - \text{FPR})) \leq R/(\eta(1 - R))\).
\(\square\)

\subsubsection{4.2 Theorem 5 Proof}\label{theorem-5-proof}

\textbf{Theorem 5} follows directly from the standard PAC-Bayes bound
(McAllester, 1999) applied to the 0-1 loss \(\ell_\theta \in [0,1]\).

\textbf{Standard PAC-Bayes bound} (McAllester, 1999, Theorem 2): For any
distribution \(\mathcal{D}\), any hypothesis class \(\mathcal{H}\), any
prior \(P\) over \(\mathcal{H}\) independent of the training data, and
any \(\delta \in (0,1)\), with probability at least \(1-\delta\) over
\(S \sim \mathcal{D}^N\), for all posteriors \(Q\) over \(\mathcal{H}\):

\[\mathbb{E}_{h \sim Q}[R(h)] \leq \mathbb{E}_{h \sim Q}[\hat{R}_S(h)] + \sqrt{\frac{\text{KL}(Q \parallel P) + \log\frac{2\sqrt{N}}{\delta}}{2N}}\]

\textbf{Application}: Take
\(\mathcal{H} = \{h_\theta : \theta \in \Theta\}\) as the detector
family, \(S = V\) as the validation set, \(R(h_\theta) = R(\theta)\) as
the true detection error, and
\(\hat{R}_V(h_\theta) = \hat{R}_V(\theta)\) as the empirical detection
error. The loss is bounded in \([0,1]\) since it is the 0-1 indicator.
Applying the standard PAC-Bayes bound yields the result. \(\square\)

\subsubsection{4.3 Corollary 5.1 Proof}\label{corollary-5.1-proof}

From Lemma 4, for any \(\theta \in \Theta\) with \(R(\theta) \leq 1/2\):

\[1 - \text{F1}(\theta) \leq \frac{2R(\theta)}{\eta}\]

Taking expectations under \(Q\):

\[\mathbb{E}_{\theta \sim Q}[1 - \text{F1}(\theta)] \leq \frac{2}{\eta} \cdot \mathbb{E}_{\theta \sim Q}[R(\theta)]\]

By Jensen's inequality (since \(1 - \text{F1}\) is convex in \(R\) for
small \(R\), but here we use linearity of expectation):

\[\mathbb{E}_{\theta \sim Q}[1 - \text{F1}(\theta)] = 1 - \mathbb{E}_{\theta \sim Q}[\text{F1}(\theta)]\]

Let \(\text{F1}(Q) = \mathbb{E}_{\theta \sim Q}[\text{F1}(\theta)]\).
Then:

\[1 - \text{F1}(Q) \leq \frac{2}{\eta} \cdot R(Q)\]

Similarly, for the empirical quantities:

\[1 - \widehat{\text{F1}}(Q) \leq \frac{2}{\eta} \cdot \hat{R}_V(Q)\]

where
\(\widehat{\text{F1}}(Q) = \mathbb{E}_{\theta \sim Q}[\widehat{\text{F1}}_V(\theta)]\)
and \(\widehat{\text{F1}}_V(\theta)\) is the empirical F1 on \(V\).

Applying Theorem 5:

\[\begin{aligned}
1 - \text{F1}(Q) &\leq \frac{2}{\eta} \cdot R(Q) \\
&\leq \frac{2}{\eta} \left[ \hat{R}_V(Q) + \sqrt{\frac{\text{KL}(Q \parallel P) + \log\frac{2\sqrt{N}}{\delta}}{2N}} \right] \\
&\leq (1 - \widehat{\text{F1}}(Q)) + \frac{2}{\eta} \sqrt{\frac{\text{KL}(Q \parallel P) + \log\frac{2\sqrt{N}}{\delta}}{2N}}
\end{aligned}\]

Rearranging:

\[\text{F1}(Q) \geq \widehat{\text{F1}}(Q) - \frac{2}{\eta} \sqrt{\frac{\text{KL}(Q \parallel P) + \log\frac{2\sqrt{N}}{\delta}}{2N}}\]

With probability \(\geq 1-\delta\). \(\square\)

\subsubsection{4.4 Corollary 5.2 Proof}\label{corollary-5.2-proof}

When \(Q = \delta_{\hat{\theta}}\) is a Dirac distribution and \(P\) is
uniform over a discretized grid of \(T\) thresholds:

\[\text{KL}(\delta_{\hat{\theta}} \parallel P) = \int \log\frac{d\delta_{\hat{\theta}}}{dP}\, d\delta_{\hat{\theta}} = \log\frac{1}{P(\hat{\theta})} = \log T\]

Substituting into Corollary 5.1 with \(|\Theta| = T\) yields the result.
\(\square\)

\begin{center}\rule{0.5\linewidth}{0.5pt}\end{center}

\subsection{5. Connection to Paper 3
Framework}\label{connection-to-paper-3-framework}

\subsubsection{5.1 Where It Fits}\label{where-it-fits}

This result corresponds to \textbf{Theorem 5} in the Paper 3 framework
(Section 4.2, ``PAC-Bayes Generalization Bound''). The framework's
novelty claim \#4 states:

\begin{quote}
\textbf{Generalized Dawid-Skene equivalence}: We prove that SCX
weighting reduces to Dawid-Skene weighting under trivial state
partition, and that the improvement is lower-bounded by the mutual
information \(I(S; W)\) between states and optimal weights.
\end{quote}

The PAC-Bayes bound extends this by providing a fully empirical
performance guarantee.

\subsubsection{5.2 How It Extends Theorem
1}\label{how-it-extends-theorem-1}

\begin{longtable}[]{@{}
  >{\raggedright\arraybackslash}p{(\linewidth - 4\tabcolsep) * \real{0.1633}}
  >{\raggedright\arraybackslash}p{(\linewidth - 4\tabcolsep) * \real{0.3878}}
  >{\raggedright\arraybackslash}p{(\linewidth - 4\tabcolsep) * \real{0.4490}}@{}}
\toprule\noalign{}
\begin{minipage}[b]{\linewidth}\raggedright
Aspect
\end{minipage} & \begin{minipage}[b]{\linewidth}\raggedright
Theorem 1 (Current)
\end{minipage} & \begin{minipage}[b]{\linewidth}\raggedright
Theorem 5 (PAC-Bayes)
\end{minipage} \\
\midrule\noalign{}
\endhead
\bottomrule\noalign{}
\endlastfoot
\(\mu_s\) knowledge & Assumed known & Estimated from \(n_s\) samples \\
Threshold selection & Requires \(\mu_s\) to set \(\theta^*\) &
Data-driven via empirical risk minimization \\
Bound type & Population & Finite-sample, high-probability \\
Dependence on \(M\) & \(\exp(-2M\Delta_s^2)\) & Captured through
\(\hat{R}_V(\theta)\) \\
Additional penalty & None &
\(\sqrt{(\text{KL}(Q\|P) + \log(1/\delta)) / (2N)}\) \\
\end{longtable}

\subsubsection{5.3 Synergy with Corollary 4 (Finite-Sample
Correction)}\label{synergy-with-corollary-4-finite-sample-correction}

The current Corollary 4 already provides a naive finite-sample
correction via Hoeffding:

\[|\hat{\mu}_s - \mu_s| \leq B\sqrt{\frac{\log(2M/\delta_0)}{2n_s}}\]

The PAC-Bayes approach improves on this in three ways: 1. \textbf{No
union bound}: The KL term automatically accounts for the complexity of
multiple states 2. \textbf{Tighter constants}: The
\(\sqrt{\text{KL}/N}\) term is often smaller than the sum of per-state
Hoeffding bounds 3. \textbf{Posterior flexibility}: We can choose \(Q\)
adaptively (e.g., concentrate on the empirically best threshold) and the
bound adapts

\subsubsection{5.4 Practical Significance}\label{practical-significance}

The PAC-Bayes bound tells SCX practitioners:

\begin{quote}
``Given \(N\) validation samples with known clean/noise labels, with
confidence \(1-\delta\), the true F1 is at least the validation F1 minus
a penalty that grows as
\(\sqrt{(\text{number of candidate thresholds evaluated}) / N}\).''
\end{quote}

This means: - For \(N = 10{,}000\), \(\delta = 0.05\),
\(|\Theta| = 100\) thresholds: penalty
\(\approx \frac{2}{\eta} \sqrt{(\log 100 + \log(2\sqrt{10000}/0.05)) / 20000} \approx \frac{2}{\eta} \cdot 0.029\)
- Even with \(\eta = 0.1\), the penalty is \(\approx 0.58\) --- why so
large? Because F1 is ill-defined near \(\eta = 0\). For practical noise
rates \(\eta \geq 0.05\), the bound becomes useful.

\begin{center}\rule{0.5\linewidth}{0.5pt}\end{center}

\subsection{6. Practical Implications}\label{practical-implications}

\subsubsection{6.1 How to Use the Bound}\label{how-to-use-the-bound}

\textbf{Step 1}: Prepare a validation set \(V\) with known clean/noise
labels (e.g., inject synthetic noise in a held-out portion of clean
data).

\textbf{Step 2}: Define a grid of candidate thresholds
\(\Theta = \{\theta_1, \dots, \theta_T\}\).

\textbf{Step 3}: For each \(\theta \in \Theta\), compute the empirical
F1 on \(V\):

\[\widehat{\text{F1}}_V(\theta) = \frac{2 \cdot \text{TP}_V(\theta)}{2 \cdot \text{TP}_V(\theta) + \text{FP}_V(\theta) + \text{FN}_V(\theta)}\]

\textbf{Step 4}: Choose the threshold \(\hat{\theta}\) that maximizes
\(\widehat{\text{F1}}_V(\theta)\).

\textbf{Step 5}: Apply Corollary 5.2 to obtain a lower bound on the true
F1:

\[\text{F1}(\hat{\theta}) \geq \widehat{\text{F1}}_V(\hat{\theta}) - \frac{2}{\eta}\sqrt{\frac{\log T + \log\frac{2\sqrt{N}}{0.05}}{2N}}\]

with 95\% confidence.

\subsubsection{6.2 When the Bound is
Tight}\label{when-the-bound-is-tight}

The PAC-Bayes bound is tight when: 1. \textbf{\(N\) is large} relative
to the complexity \(\text{KL}(Q\|P)\) 2. \textbf{The empirical F1 is
accurate} (validation set is representative of the true distribution) 3.
\textbf{\(R(\theta) \ll 1/2\)} so the F1-risk conversion is efficient

\subsubsection{6.3 When the Bound is
Loose}\label{when-the-bound-is-loose}

The bound becomes loose when: 1. \textbf{\(N\) is small} (\(< 100\)
samples per state): the complexity penalty dominates 2. \textbf{\(\eta\)
is very small} (\(< 0.01\)): the \(1/\eta\) amplification makes the
bound vacuous 3. \textbf{The empirical F1 is near 0 or 1}: the 0-1 loss
PAC-Bayes bound doesn't leverage the variance at the extremes (use the
empirical Bernstein variant in Item 3 instead)

\subsubsection{6.4 Relation to Catoni's
PAC-Bayes}\label{relation-to-catonis-pac-bayes}

Catoni (2007) provides a refined PAC-Bayes bound with sharper constants.
The bound above uses McAllester's original formulation for simplicity.
In the actual paper, we can use Catoni's bound:

\[R(Q) \leq \hat{R}_V(Q) + \frac{\text{KL}(Q \parallel P) + \log(1/\delta)}{\lambda N} + \frac{\lambda}{8N}\]

where \(\lambda > 0\) is an inverse-temperature parameter. Optimizing
over \(\lambda\) gives a tighter bound, especially when \(\hat{R}_V(Q)\)
is small.

\begin{center}\rule{0.5\linewidth}{0.5pt}\end{center}

\subsection{7. Open Questions and
Extensions}\label{open-questions-and-extensions}

\begin{enumerate}
\def\labelenumi{\arabic{enumi}.}
\item
  \textbf{Data-dependent prior}: Can we use a prior that depends on the
  structure of the state partition? The current bound requires \(P\) to
  be independent of \(V\).
\item
  \textbf{PAC-Bayes with F1 directly}: The standard PAC-Bayes bound
  applies to the 0-1 loss, not to the F1 score directly. A PAC-Bayes
  bound for the F1 itself would be tighter but requires more
  sophisticated tools (e.g., the F1's difference-of-convex structure).
\item
  \textbf{Unlabeled validation set}: In the pure SCX setting, we may not
  have ground-truth labels for the validation set. Can we bound the
  detector's performance using only clean validation samples? This would
  require a different approach --- perhaps using the estimated
  \(\hat{\mu}_s\) directly and propagating uncertainty.
\item
  \textbf{Semi-supervised PAC-Bayes}: Extend to settings where only a
  subset of the validation set has known labels (active learning for
  threshold calibration).
\end{enumerate}

\begin{center}\rule{0.5\linewidth}{0.5pt}\end{center}

\subsection{References}\label{references}

\begin{enumerate}
\def\labelenumi{\arabic{enumi}.}
\item
  McAllester, D. A. (1999). PAC-Bayesian model averaging.
  \emph{Proceedings of the Twelfth Annual Conference on Computational
  Learning Theory}.
\item
  Catoni, O. (2007). PAC-Bayesian supervised classification: The
  thermodynamics of statistical learning. \emph{IMS Lecture Notes
  Monograph Series}, 56.
\item
  Germain, P., Bach, F., Lacoste, A., \& Marchand, M. (2009).
  PAC-Bayesian learning of linear classifiers. \emph{Proceedings of
  ICML}.
\item
  Theorem 1 (SCX Noise Detection).
  \texttt{../theorems/01\_noise\_detection\_guarantee.md}
\item
  Corollary 4 (Finite-Sample Correction).
  \texttt{../theorems/01\_noise\_detection\_guarantee.md\#34-\%E6\%9C\%89\%E9\%99\%90\%E6\%A0\%B7\%E6\%9C\%AC\%E6\%A0\%A1\%E6\%AD\%A3-finite-sample-correction}
\end{enumerate}

\end{document}
